% Lösungen Extremalprobleme · Maximum bestimmen · Prüfungs-Fokus Abitur GK (zusammenbau v0.6)
\einheitenkopf{Einheit 3 · Maximum bestimmen und deuten}

\erg{1}{a) der Flächeninhalt \quad b) $A'(a) = 30 - 4a = 0$ liefert $a = 7{,}5$ (im Definitionsbereich) \quad c) $A'(a) = 24 - 4a = 0$ liefert $a = 6$; $A''(6) = -4 < 0$, also Maximum; $b = 24 - 12 = 12$ m, $a = 6$ m, $A(6) = 72$ m²}
\erg{2}{a) $A(u) = u \cdot f(u) = (u^2 + 3u) \cdot e^{-0{,}5u}$; $A'(u) = 0$ liefert $u^2 - u - 6 = 0$, also $u = 3$ ($u = -2$ entfällt); $P(3 | 1{,}34)$ \quad b) $A(u) = (u^2 + 4u) \cdot e^{-u}$; $A'(u) = 0$ liefert $u^2 + 2u - 4 = 0$, also $u = -1 + \sqrt{5} \approx 1{,}24$ (die negative Lösung entfällt); $P(1{,}24 | 1{,}52)$ \quad c) $A(u) = u \cdot (9 - u^2) = 9u - u^3$; $A'(u) = 9 - 3u^2 = 0$ liefert $u = \sqrt{3} \approx 1{,}73$; $A''(u) = -6u < 0$; $P(1{,}73 | 6)$ \quad d) $d(x) = f(x) - g(x) = -0{,}01x^2 + 0{,}5x$; $d'(x) = -0{,}02x + 0{,}5 = 0$ liefert $x = 25$; $d(25) = -6{,}25 + 12{,}5 = 6{,}25$ m \quad e) $d(x) = -0{,}1x^3 + 0{,}3x^2 + 0{,}9x$; $d'(x) = -0{,}3x^2 + 0{,}6x + 0{,}9 = 0$ liefert $x = 3$ oder $x = -1$ (außerhalb); Randwerte $d(0) = 0$ und $d(4{,}5) \approx 1{,}01$; größte Dicke $d(3) = 2{,}7$ dm, kleiner als $2{,}8$ dm – eingehalten \quad f) $d(x) = p(x) - f(x) = x^2 - x^4$; $d'(x) = 2x - 4x^3 = 0$ liefert $x = 0$ (dort ist $d = 0$, ein Minimum) und $x^2 = \tfrac{1}{2}$; dort $d = \tfrac{1}{2} - \tfrac{1}{4} = \tfrac{1}{4}$ \quad g) $T(u) = (10 - u) \cdot (10u - u^2) = u \cdot (10 - u)^2$; $T'(u) = (10 - u) \cdot (10 - 3u) = 0$ liefert $u = 10$ (nicht im Bereich) oder $u = \tfrac{10}{3} \approx 3{,}33$; Vorzeichenwechsel von $+$ nach $-$ – Maximum \quad h) $T(u) = 0{,}5u \cdot (6 - u)^2$; $T'(u) = 0{,}5 \cdot (6 - u) \cdot (6 - 3u) = 0$ liefert $u = 6$ (nicht im Bereich) oder $u = 2$; Vorzeichenwechsel von $+$ nach $-$; $T(2) = 16$ \quad i) $T(u) = 2u \cdot (4 - u)^2$; $T'(u) = 2 \cdot (4 - u) \cdot (4 - 3u) = 0$ liefert $u = 4$ (nicht im Bereich) oder $u = \tfrac{4}{3} \approx 1{,}33$; Vorzeichenwechsel von $+$ nach $-$ – Maximum \quad j) $A(x) = x \cdot f(x) = 4x - x^3$, $A'(1) = 4 - 3 = 1 \ne 0$ – die notwendige Bedingung ist verletzt, also ist $x = 1$ keine Maximalstelle. \quad k) $A(x) = 6x - 2x^2$, $A'(2) = 6 - 8 = -2 \ne 0$ – die notwendige Bedingung ist verletzt, $x = 2$ ist keine Maximalstelle. \quad l) $A(x) = x^4 - 5x^3 + 6x^2 + x$, $A'(x) = 4x^3 - 15x^2 + 12x + 1$, $A'(1) = 4 - 15 + 12 + 1 = 2 \ne 0$ – keine Maximalstelle (nicht $f'(1)$ prüfen).}
